
\begin{abstract}
Komunikasi antaragen dalam \textit{Multi-Agent System} (MAS) berbasis 
\textit{Large Language Model} (LLM) umumnya menggunakan teks, 
sedangkan komunikasi berbasis ruang laten berpotensi mengurangi kebutuhan 
representasi tekstual. Namun, bagaimana formulasi penalaran laten dan 
medium komunikasi memengaruhi kinerja MAS masih perlu dikaji. Penalaran laten 
memungkinkan LLM memproses informasi melalui representasi internal, 
sementara komunikasi laten melalui \textit{key-value cache} (KV-cache) membuka 
peluang bagi agen untuk bertukar informasi tanpa harus mengubahnya 
menjadi teks terlebih dahulu. \textit{Mixture of Inputs} (MoI) menggunakan 
pendekatan Bayesian untuk menggabungkan beberapa kemungkinan representasi 
berdasarkan bobot posterior, sehingga diharapkan dapat menghasilkan masukan 
yang lebih adaptif terhadap informasi yang tersedia dan mendukung proses penalaran berikutnya. 
Penelitian ini membandingkan empat formulasi penalaran laten, 
yaitu \textit{raw}, \textit{soft}, \textit{sample}, dan MoI, 
serta dua medium komunikasi, yaitu teks dan KV-cache, menggunakan 
model \textit{open-weight} Qwen3-8B. Pengujian dilakukan pada 
\textit{benchmark} GSM8K, ARC-Challenge, dan HumanEval+ dengan mengevaluasi 
akurasi penyelesaian tugas, penggunaan token, dan waktu inferensi. 
Perbandingan ini diharapkan dapat menunjukkan sejauh mana komunikasi 
berbasis KV-cache meningkatkan efisiensi dibandingkan komunikasi berbasis teks, 
sekaligus mengungkap pengaruh formulasi penalaran laten terhadap keseimbangan 
antara efisiensi dan akurasi.

\textbf{Kata kunci:} \textit{Mixture of Inputs}, penalaran laten, 
\textit{Multi-Agent System}, KV-cache, \textit{Large Language Model}.
\end{abstract}
